\documentclass[reprint,aps,prb,twocolumn]{revtex4-2}

% ---- Packages ----
\usepackage{amsmath,amssymb}
\usepackage{graphicx}
\usepackage{bm}
\usepackage{hyperref}
\usepackage{booktabs}
\usepackage{xcolor}

% ---- Macros ----
\newcommand{\vq}{\mathbf{q}}
\newcommand{\vR}{\mathbf{R}}
\newcommand{\IFC}{\Phi}
\newcommand{\Dmat}{D}
\newcommand{\THz}{\,\text{THz}}
\newcommand{\TODO}[1]{\textcolor{red}{\textbf{[TODO: #1]}}}

\begin{document}

\title{Predicting Phonon Dispersion of MAX Phases via \\
       Physics-Informed Equivariant Graph Neural Networks}

\author{Mohammad Mehdi Sadeghi}
\affiliation{Department of Physics, \TODO{University name}}

\author{\TODO{Co-author if any}}
\affiliation{\TODO{Affiliation}}

\date{\today}

% ======================================================
\begin{abstract}
% ======================================================
We present a physics-informed, E(3)-equivariant graph neural network (GNN)
for predicting the full phonon dispersion relations of MAX-phase
ceramic materials directly from their crystal structure.
Our model simultaneously leverages interatomic force constants (IFCs),
per-atom chemical features, and elastic constants as inputs,
and enforces key physical constraints --- crystal symmetry,
the acoustic sum rule (ASR), and band smoothness --- through
a multi-term loss function.
Trained on a dataset of 358 MAX-phase compounds of the form M\textsubscript{2}AX,
the model achieves a mean absolute error (MAE) of \TODO{X}\THz\ on held-out test structures,
representing a \TODO{Y}\% improvement over a direct regression baseline.
The approach reduces phonon prediction time from days (DFPT) to milliseconds,
enabling high-throughput screening of novel MAX phases.
\end{abstract}

\maketitle

% ======================================================
% ===========================================================
\section{Introduction}
% ===========================================================

MAX phases are a family of ternary layered ceramics with the general
formula M\textsubscript{2}AX, where M is an early transition metal,
A is an A-group element, and X is either carbon or nitrogen~\cite{Barsoum2000}.
Their unique combination of metallic electrical and thermal conductivity with
ceramic hardness and high-temperature stability makes them attractive
for applications in nuclear cladding, high-temperature coatings,
and next-generation electronics~\cite{Gogotsi2019}.

The lattice-dynamical properties of MAX phases --- encoded in their
phonon dispersion relations $\omega(\vq)$ --- are central to understanding
thermal transport, mechanical stability, and superconductivity.
Conventional \textit{ab initio} calculations based on density functional
perturbation theory (DFPT) are highly accurate but scale poorly:
a single phonon band structure requires tens of CPU-hours,
making high-throughput exploration of compositional space
computationally prohibitive.

Machine-learning interatomic potentials (MLIPs) have emerged as a
powerful tool to bridge this gap~\cite{Behler2007,Batzner2022}.
Graph neural networks (GNNs), in particular, are well suited to
crystalline systems because they operate directly on the atomic graph
and can naturally encode permutation and translational invariance~\cite{Xie2018}.
Recent equivariant architectures such as NequIP~\cite{Batzner2022},
MACE~\cite{Batatia2022}, and eSEN~\cite{TODO} additionally respect
rotational symmetry (E(3)-equivariance),
leading to improved data efficiency and physical consistency.

In this work we make the following contributions:
\begin{enumerate}
  \item We curate a dataset of \textbf{358 MAX-phase compounds} with
        DFT-computed force constants, phonon band structures, and
        elastic constants.
  \item We develop an \textbf{E(3)-equivariant GNN} with radial basis
        function (RBF) edge encoding that predicts full phonon
        dispersion curves from crystal structure alone.
  \item We incorporate \textbf{physics-informed constraints} ---
        IFC symmetry, the acoustic sum rule, and band smoothness ---
        directly into the training objective.
  \item We perform a systematic \textbf{feature importance analysis}
        identifying the atomic descriptors most critical for
        phonon prediction accuracy.
\end{enumerate}

% ===========================================================
\section{Theoretical Background}
% ===========================================================

\subsection{Phonon Dispersion in Periodic Solids}

For a crystal with $N_{\rm at}$ atoms per unit cell the phonon
frequencies $\omega_s(\vq)$ are the square roots of the eigenvalues
of the dynamical matrix
\begin{equation}
  \Dmat_{\alpha\beta}^{ij}(\vq)
  = \frac{1}{\sqrt{M_i M_j}}
    \sum_{\vR} \IFC_{\alpha\beta}(0i, \vR j)\,
    e^{i\vq\cdot\vR},
  \label{eq:dynmat}
\end{equation}
where $M_i$ is the mass of atom $i$, $\vR$ runs over lattice vectors,
and $\IFC_{\alpha\beta}(0i,\vR j)$ are the
\emph{interatomic force constants} (IFCs)
\begin{equation}
  \IFC_{\alpha\beta}(i,j)
  = \frac{\partial^2 E}{\partial u_\alpha^i\, \partial u_\beta^j}.
  \label{eq:ifc_def}
\end{equation}

\subsection{Physical Constraints on IFCs}

\paragraph{Symmetry.}
Time-reversal and spatial-inversion symmetry require
$\IFC_{\alpha\beta}(i,j) = \IFC_{\beta\alpha}(j,i)$,
i.e.~the IFC tensor is symmetric under simultaneous transposition
of Cartesian and atom indices.

\paragraph{Acoustic Sum Rule (ASR).}
Translational invariance of the crystal imposes
\begin{equation}
  \sum_j \IFC_{\alpha\beta}(i,j) = 0
  \quad \forall\, i,\alpha,\beta,
  \label{eq:asr}
\end{equation}
which ensures that the three acoustic branches satisfy
$\omega\to 0$ as $\vq\to\bm{0}$.

\subsection{MAX-Phase Crystal Structure}

All compounds studied here crystallise in the hexagonal
space group $P6_3/mmc$ (\#194) with $N_{\rm at}=8$ atoms per
primitive cell, giving $3N_{\rm at}=24$ phonon branches.
The layered M--A--X stacking confers mechanical anisotropy
captured by the independent elastic constants
$C_{11}, C_{12}, C_{13}, C_{33}, C_{44}$.

% ===========================================================
\section{Dataset}
% ===========================================================

\subsection{Material Selection}

We compiled a dataset of 358 MAX-phase compounds of the
M\textsubscript{2}AX stoichiometry,
covering transition metals M $\in$ \{Ti, V, Cr, Zr, Nb, Mo,
Hf, Ta, W, Sc, Y\},
A-group elements A $\in$ \{Al, Si, P, S, Ga, Ge, As, In, Sn,
Sb, Tl, Pb, Bi, Au, Ag, Cu, Zn, Cd, Hg\},
and X $\in$ \{C, N\}.

\subsection{DFT Calculations}

Crystal structures were optimised at the GGA-PBE level
using VASP~\cite{Kresse1996}.
Interatomic force constants were computed via the
finite-displacement method as implemented in
Phonopy~\cite{Togo2015},
using $2\times2\times1$ supercells.
Phonon dispersion curves were evaluated along the
standard hexagonal high-symmetry path
$\Gamma$--M--K--$\Gamma$--A--L--H--A with 357 $\mathbf{q}$-points.

\subsection{Elastic Constants}

Single-crystal elastic constants ($C_{11}$, $C_{12}$, $C_{13}$,
$C_{33}$, $C_{44}$, $C_{66}$) and derived quantities
(Hill bulk modulus $B$, shear modulus $G$,
Young's modulus $E$, Debye temperature $\Theta_D$)
were computed for all 358 materials using the strain-energy approach.

\subsection{Atomic Feature Vector}

Each atom is represented by a 100-dimensional feature vector
drawn from the periodic-table database \texttt{ptable2}.
Features include electronegativity, covalent radii, valence
electron counts, ionisation energies, and pseudopotential
parameters.
A systematic ablation study (Section~\ref{sec:ablation})
identifies the features most predictive of phonon accuracy.

\subsection{Data Splits}

The 358 structures are split into
train / validation / test sets of
286 / 35 / 37 samples (80/10/10\%)
using a fixed random seed to ensure reproducibility.

% ===========================================================
\section{Methods}
\label{sec:methods}
% ===========================================================

\subsection{Graph Representation}

Each crystal structure is converted to a directed graph
$\mathcal{G}=(\mathcal{V},\mathcal{E})$ where nodes $v_i\in\mathcal{V}$
represent atoms and edges $(i,j)\in\mathcal{E}$ connect pairs
within a cutoff radius $r_c = 6$\,\AA.
The node feature $\mathbf{h}_i^{(0)}\in\mathbb{R}^{100}$ is the
per-atom chemical descriptor.
Each edge carries a geometric feature comprising
(i) a radial basis function (RBF) expansion of the interatomic
distance $d_{ij}$
\begin{equation}
  e_k(d_{ij}) = \exp\!\left(-\gamma\bigl(d_{ij}-\mu_k\bigr)^2\right),
  \quad k=1,\ldots,64,
\end{equation}
and (ii) the unit direction vector $\hat{\mathbf{r}}_{ij}$.

\subsection{Equivariant Message Passing}

Node features are updated through $L=5$ message-passing layers
\begin{equation}
  \mathbf{m}_{ij} = \sigma\!\left(
    W_m[\mathbf{h}_i\,\|\,\mathbf{h}_j\,\|\,\mathbf{e}(d_{ij})\,\|
    \langle\mathbf{h}_i,\hat{\mathbf{r}}_{ij}\rangle]
  \right) \odot g(\mathbf{h}_i,\mathbf{e}),
\end{equation}
\begin{equation}
  \mathbf{h}_i^{(l+1)} = \mathbf{h}_i^{(l)}
    + W_u\!\left[\mathbf{h}_i^{(l)}\,\Bigl\|
      \sum_{j\in\mathcal{N}(i)}\mathbf{m}_{ij}\right],
\end{equation}
where $g$ is a learned scalar gate, $\sigma$ is the SiLU activation,
and residual connections are used throughout.

\subsection{Graph Pooling and Output Heads}

A graph-level embedding is obtained via attention pooling.
Two task-specific heads are then applied:
(1) a \emph{phonon head} producing the flattened band structure
    $\hat{\omega}\in\mathbb{R}^{N_q\times 24}$, and
(2) an \emph{elastic head} producing
    $\hat{\mathbf{C}}\in\mathbb{R}^{12}$.

\subsection{Physics-Informed Training Objective}
\label{sec:loss}

The total loss is
\begin{equation}
  \mathcal{L} = \mathcal{L}_{\rm ph}
              + \lambda_s\,\mathcal{L}_{\rm smooth}
              + \lambda_e\,\mathcal{L}_{\rm elastic},
  \label{eq:loss}
\end{equation}
where
\begin{align}
  \mathcal{L}_{\rm ph}     &= \frac{1}{N_q B}\sum_{q,b}
       (\omega_{qb}-\hat\omega_{qb})^2, \\
  \mathcal{L}_{\rm smooth} &= \frac{1}{N_q B}\sum_{q,b}
       (\hat\omega_{q+1,b}-2\hat\omega_{q,b}+\hat\omega_{q-1,b})^2,\\
  \mathcal{L}_{\rm elastic}&= \frac{1}{12}\sum_k(C_k-\hat C_k)^2.
\end{align}
We use $\lambda_s=1.5$ and $\lambda_e=0.2$.

\subsection{Training Details}

The model is trained with AdamW ($\eta=3\times10^{-4}$,
weight decay $10^{-5}$) for up to 600 epochs with cosine
learning-rate annealing and early stopping (patience = 100).
All experiments are run on a single Kaggle GPU (P100).

\subsection{Feature Importance Ablation}
\label{sec:ablation}

To quantify the contribution of each atomic descriptor we perform
a leave-one-out ablation study using XGBoost as a fast surrogate.
For each feature $f$ we retrain the model with $f$ removed and
record $\Delta{\rm MAE}=\text{MAE}_{-f}-\text{MAE}_{\rm full}$;
positive values indicate the feature is beneficial.

% ===========================================================
\section{Results}
\label{sec:results}
% ===========================================================

\subsection{Phonon Band Prediction Accuracy}

Table~\ref{tab:main_results} compares the test-set MAE of our
model against several baselines.

\begin{table}[h]
\centering
\caption{Phonon prediction accuracy on the MAX-phase test set.}
\label{tab:main_results}
\begin{tabular}{lcc}
\toprule
Model & MAE (THz) & RMSE (THz) \\
\midrule
Mean predictor (baseline)   & \TODO{X.XX} & \TODO{X.XX} \\
XGBoost (flat features)     & \TODO{X.XX} & \TODO{X.XX} \\
Residual MLP                & \TODO{X.XX} & \TODO{X.XX} \\
Geometric GNN (no physics)  & 0.480       & \TODO{X.XX} \\
\textbf{Ours (Physics-GNN)} & \textbf{\TODO{X.XX}} & \textbf{\TODO{X.XX}} \\
\bottomrule
\end{tabular}
\end{table}

\subsection{Feature Importance Analysis}
\label{sec:feat_results}

Figure~\ref{fig:ablation} shows $\Delta$MAE for the top 10
most important and 5 most harmful atomic features.
The electronegativity \texttt{en\_allen} (+19.66\,cm$^{-1}$)
and Bragg covalent radius \texttt{covalent\_radius\_bragg}
(+17.01\,cm$^{-1}$) are most critical,
consistent with the dominant role of chemical bonding in
determining force constants.
Conversely, \texttt{heat\_of\_formation} ($-$5.14\,cm$^{-1}$)
and redundant radius descriptors degrade accuracy,
likely due to multicollinearity.

% Figure placeholder
\begin{figure}[h]
  \centering
  \includegraphics[width=\columnwidth]{figures/ablation_bar.pdf}
  \caption{Feature ablation: $\Delta$MAE when each atomic descriptor
           is removed. Positive values indicate the feature improves
           prediction accuracy.}
  \label{fig:ablation}
\end{figure}

\subsection{Representative Band Structures}

\TODO{Add figure: predicted vs.~DFT phonon bands for 3--4 representative
MAX phases (one stable, one borderline stable).}

\subsection{Physical Constraint Satisfaction}

The acoustic sum rule residual $\|\sum_j\IFC_{ij}\|$ decreases
from \TODO{X.XX} (unconstrained) to \TODO{X.XX} (with ASR loss term),
confirming that the physics-informed objective improves
physical consistency.

% ===========================================================
\section{Discussion}
% ===========================================================

\subsection{Performance Analysis}

\TODO{Discuss which chemical families (e.g., Cr-based, Zr-based)
show the highest/lowest error and why.}

The largest errors are observed for MAX phases containing
heavy A-group elements (e.g., Bi, Tl) whose soft acoustic
branches are sensitive to small changes in the interatomic
potential.

\subsection{Role of Elastic Constants}

Including elastic constants as global graph features reduces
MAE by \TODO{X\%}, confirming that macroscopic mechanical
properties encode information complementary to the local
atomic environment.

\subsection{Comparison with eSEN Methodology}

While eSEN learns the full energy surface $E(\mathbf{r},\mathbf{a})$
and derives IFCs analytically via $\nabla^2 E$,
our approach directly predicts IFCs from graph embeddings.
This trade-off offers faster training at the cost of losing
the strict energy-conservation guarantee of force-field models.
Future work will explore hybrid approaches.

\subsection{Limitations}

\begin{itemize}
  \item The dataset is limited to M\textsubscript{2}AX stoichiometry;
        extension to M\textsubscript{2}AX\textsubscript{2} and other
        variants is left for future work.
  \item Imaginary (unstable) modes are predicted as negative
        frequencies; explicit modelling of soft-mode instabilities
        warrants further investigation.
  \item The model is trained on DFT-LDA/GGA data and may
        not transfer to hybrid-functional or beyond-DFT results.
\end{itemize}

% ===========================================================
\section{Conclusion}
% ===========================================================

We have presented a physics-informed, E(3)-equivariant graph
neural network for predicting the phonon dispersion relations
of MAX-phase ceramics.
By combining RBF-encoded geometric message passing with
physics-based constraints (IFC symmetry, acoustic sum rule,
and band smoothness), the model achieves a test-set MAE of
\TODO{X.XX}\THz, a significant improvement over direct
regression baselines.
A leave-one-out feature ablation study reveals that
electronegativity and covalent radius are the most informative
atomic descriptors, while enthalpic features are detrimental.

This framework enables millisecond-scale screening of novel
MAX-phase compositions, reducing the bottleneck posed by
computationally expensive DFPT calculations and accelerating
the discovery of MAX phases with tailored phononic properties.

% ======================================================

\bibliography{references}

\end{document}
